\section{Conclusion}\label{sec:conclusion}

This paper asks how spatial price competition changes when destination demand reallocates a scarce transport-service resource. The key feedback is two-sided: a retailer's price affects shopping demand directly, and that demand shifts service toward its own direction and away from its rival. Directional background demand makes the feedback asymmetric.

The model shows that this mechanism can support a nonempty open set of global price equilibria in which one retailer treats the rival's price as a strategic substitute while the other retains the conventional complementary response. A minimum service obligation can remain slack at that equilibrium yet matter for global support by truncating extreme off-path service withdrawal. The transport allocation problem also provides a coherent real-cost benchmark, under which decentralized pricing generally differs from the constrained-efficient allocation subject to the same service floor.

The interpretation is deliberately narrow. Strategic asymmetry, network effects, demand--frequency feedback, and minimum transit-service standards are established ideas; the contribution is their full-game interaction through a third-party shared service resource. The model is most natural for large destination zones, managed directional or shuttle services, and demand-responsive or shared-mobility systems. It does not derive an optimal service standard, subsidy, fleet size, or retailer location.
